\chapter{总结与展望}
\label{chp:conclusion}

\section{研究总结}
\label{sec:conclusion}

随着人工智能与机器人技术的不断发展，具身智能作为连接数字世界与物理世界的桥梁，正成为科技发展的重点领域。然而，真实世界中的非结构化环境往往伴随着光照变化、物体层叠遮挡等复杂因素，同时长时序任务中还存在明显的历史上下文依赖，这些问题共同影响了机械臂在复杂操作任务中的成功率与可靠性。针对长时序、非结构化任务中的多模态动作生成与决策依赖难题，本文提出了一种基于特征时序缓存的视觉语言动作模型 TemporalCacheVLA。通过在仿真数据集与真实机械臂平台上的实验评估，结果表明，本文所提出的方法在多阶段、长时序操作任务中具有较好的性能表现，并在一定程度上验证了其真实部署潜力。本文的主要研究工作与结论如下：

（1）研究了基于运动学理论与典型模仿学习算法的机器人操作基础理论。针对机械臂在笛卡尔空间与关节空间之间的运动映射问题，本文基于 Denavit-Hartenberg 参数法对 Franka Panda 与 UR5e 两款机械臂进行了运动学建模。该模型推导了正向与逆向运动学方程，明确了各关节角度与末端空间位姿之间的映射关系，为后续真实场景数据集构建中的动作标签计算提供了正运动学支撑，同时也为真实部署阶段末端增量动作到关节空间控制目标的转换提供了理论基础。

（2）研究了基于扩散模型的模仿学习动作生成算法。针对传统行为克隆方法难以有效拟合多模态动作分布的问题，提出了一种结合多模态特征编码与扩散去噪的动作生成算法。该方法采用迭代 DiT 为核心架构，利用 EfficientViT 提取轻量化多尺度视觉特征。为了实现多模态特征的有效融合，设计了非对称跨模态注意力机制，在保留状态残差的同时引导多步去噪过程，并采用 DDIM 策略加速采样。实验结果表明，在 RoboMimic 数据集上，该算法在多项任务中取得了较好的性能表现。相比于基线模型，本文方法在 Transport ph 任务上取得了 $0.862$ 的成功率，在 Tool Hang ph 任务上达到 $0.791$；此外，在 Transport mh 数据集上的对比实验中，本文方法的成功率由 Diffusion Policy 的 $0.643$ 提高至 $0.801$。这些结果说明，所建立的扩散动作生成模型能够较有效地建模复杂多模态动作分布，并在生成质量与推理效率之间取得较好的平衡。

（3）研究了基于特征时序缓存的视觉语言动作模型 TemporalCacheVLA。针对机器人在长时序任务中因单帧决策导致的状态混淆与上下文遗忘问题，本文提出了一种双通道特征时序缓存机制，并在此基础上构建了 TemporalCacheVLA 模型。该模型通过感知语义双通道历史存储池分别维护历史感知信息与历史语义信息，并结合缓存检索、可学习门控融合以及基于余弦相似度的压缩更新策略，实现了对长时序上下文的高效维护与动态利用。在具体实现上，模型采用双分支视觉编码器与大语言模型进行多模态特征提取，并将增强后的上下文特征输入扩散动作生成模块。实验结果表明，在 LIBERO 基准数据集上，本文方法在强调指令跟随的 Goal 子集与包含长时序任务的 LIBERO-100 子集上分别取得了 $0.959$ 和 $0.949$ 的平均成功率。同时，在 LIBERO-100 子集上的消融实验中，移除缓存检索、门控融合与缓存压缩机制后，成功率分别下降至 $0.862$、$0.835$ 和 $0.823$。这些结果说明，双通道特征时序缓存机制能够较有效地维护和利用历史上下文信息，从而缓解长时序任务中的决策波动问题。

（4）研究了基于真实物理平台的系统部署与实验验证。针对视觉语言动作模型向真实物理世界迁移部署时面临的数据采集与系统集成难题，本文搭建了基于 GELLO 遥操臂的遥操作数据采集平台。该平台利用与 UR5e 运动学等价的缩比遥操臂采集包含纠错轨迹的示教数据，对模型进行参数冻结微调，并基于 ROS 框架开发了包含会话校验与动作限幅的本地控制节点。在真实机械臂平台上的对比实验表明，TemporalCacheVLA 在简单抓取、抽屉操作及长时序组合任务中分别取得了 $0.85$、$0.65$ 和 $0.40$ 的成功率。同时，相比于 OpenVLA 基线模型，本文方法在长时序组合任务中的成功率提升了 $0.30$。这些结果初步表明，本文所提出的框架在真实物理环境中具有一定的应用可行性，并展现出较好的长时序任务建模潜力。

\section{研究展望}
\label{sec:future_work}

本文围绕基于扩散模型与时序缓存的视觉语言动作机械臂自主操作技术开展了较为系统的研究，并完成了仿真与真实平台上的初步实验评估。然而，受限于现有硬件条件与数据获取方式，本研究仍存在一些不足。首先，模型参数量较大，推理开销较高，难以在算力受限的边缘设备上实现单机部署；其次，依赖遥操臂采集的示教数据往往包含一定人为噪声，且采集成本较高，数据质量与多样性受限，进而影响了真机推理时动作的平滑性与泛化能力；此外，当前模型在面对新任务时仍需重新采集数据并进行微调，尚不具备开箱即用的通用泛化能力。针对上述不足，未来还可从以下几个方向进行拓展与深入：

（1）探索面向边缘部署的模型轻量化技术。目前的大型视觉语言动作模型在实际部署时，仍对计算资源有较高要求，难以在算力受限的边缘设备上实现低延迟推理。未来可以引入模型剪枝、知识蒸馏与低比特量化等轻量化技术，在保证模型语义理解与动作生成精度的前提下，大幅压缩模型体积与计算复杂度，从而推动机械臂自主操作系统的在线部署与实用化发展。

（2）扩展多源异构的动作学习来源。当前的模仿学习主要依赖于遥操臂采集的专家示教数据，数据获取成本较高，且覆盖的操作模式有限，容易导致模型在面对全新物体或场景时泛化受限。未来可以引入基于手势分析模型的人类视频演示学习技术，结合多源感知手段，直接从海量互联网人类操作视频中提取动作先验，以丰富可学习的动作库，从而进一步增强模型在开放场景中的适应性。

（3）结合在线视觉语言大模型提升跨任务通用性与泛化能力。现有的视觉语言动作模型通常需要针对特定任务重新采集数据并进行微调，这限制了系统在面对未见任务、未见物体和新场景时的泛化能力。未来可以探索将先进的在线视觉语言大模型与机器人控制系统进一步结合，利用其更强的零样本语义理解能力、跨任务知识迁移能力和开放场景推理能力，降低对任务特定示教数据与重复微调的依赖，从而提升模型在复杂开放环境中的通用性与适应性。

